\daytitle{1}{初高中物理的区别与学习方法}

\begin{learningbox}{本节导读}
初中物理已经包含定性分析和定量计算；进入高中后，问题会更强调理想模型、矢量方向、公式的适用条件以及多步骤推理。本节先建立高中物理的基本框架：什么是模型？什么是矢量？为什么单位和量纲重要？拿到一道高中物理题应当怎样组织思路？
\end{learningbox}

\textbf{本节读题路线：} 先把题目中的真实情境翻译成研究对象和理想模型，再判断涉及的是矢量还是标量；写公式前统一单位，写答案后用量纲和数量级检查。本节先不追求难题，而是建立“模型、方向、单位”这三个基本意识。

\section*{一、学习目标}

\begin{enumerate}
  \item 能说出高中物理相对初中物理新增的三类要求：模型、矢量和定量推理。
  \item 能用质点模型举例说明：为什么物理会暂时忽略形状、大小等次要因素。
  \item 能说出质点模型保留了什么、忽略了什么。
  \item 能区分矢量和标量，并说明方向改变会怎样影响运算结果。
  \item 能用国际单位制（SI）统一单位，并用量纲做初步检查。
  \item 能按“建模、定方向、统一单位、检查量纲”的流程处理入门题。
\end{enumerate}

\section*{二、初高中衔接}

初中物理已经学习了“力是改变物体运动状态的原因”，也进行过速度、密度、功率等定量计算。高中物理会在此基础上进一步追问：研究对象怎样简化？方向怎样规定？公式在什么条件下成立？多个规律怎样连接起来？

举个例子：已知“一个物体受到 $10\ \mathrm N$ 的力”，还不足以确定其运动变化。若进一步给出物体质量为 $2\ \mathrm{kg}$、所受合力水平向右，并采用质点模型，才可由后面将学习的牛顿第二定律得到大小为 $5\ \mathrm{m/s^2}$、方向向右的加速度。

\textbf{直观概括：} 高中物理不是简单地把计算变难，而是更重视“模型怎样建立、方向怎样处理、条件是否满足、推理能否闭合”。

\textbf{科学发展视角：} 现代科学和工程并不是直接面对完整世界，而是先建模再计算。天气预报、自动驾驶、航天轨道、芯片散热都要先确定研究对象、忽略次要因素、检查单位和数量级。读这类真实问题时，先问“保留哪些量、忽略哪些量、结果数量级是否合理”，本节学习的模型、矢量和量纲，就是以后读懂真实科技问题的底层语言。

\section*{三、核心概念}

\begin{conceptbox}{物理模型}

\textbf{直观理解：} 物理模型就是“简化版的世界”。现实中的物体有形状、颜色、温度、弹性、粗糙程度……但在物理题里，只保留跟问题相关的少数几个性质，其他因素暂时忽略。

\textbf{重要说明：} 高中物理中的质点、光滑面、轻绳、轻杆等，都不是现实中完美存在的实际对象，而是为了抓住主要矛盾建立的\textbf{理想物理模型}。

\begin{itemize}
  \item 研究苹果的平动——可把苹果看作一个有质量但没有大小的\textbf{质点}，忽略其大小、形状和转动；若题目还要求忽略空气阻力，则这是另一项独立的理想化条件。
  \item 研究桌子的受力——桌子看做不会变形的\textbf{刚体}，忽略微小形变。
  \item 研究光滑斜面——把接触面间的摩擦力取为零。这不是说现实摩擦一定很小，而是题设指定的理想化条件；若题目给出粗糙面，就必须把摩擦力重新纳入模型。
\end{itemize}

\textbf{严格定义：} 物理模型是对实际物理系统的一种抽象和简化，保留研究对象的核心物理属性，忽略次要因素，以便用数学语言精确描述其行为。

\textbf{为什么要用模型？} 因为现实太复杂。如果把所有因素都同时考虑进去，一个看似普通的苹果下落问题也会牵涉形状、空气阻力、转动、弹性等大量因素，反而看不清主导规律。物理学的智慧在于：知道忽略什么，比知道保留什么更重要。

\textbf{模型判断三问：}
\begin{enumerate}
  \item 这个因素会不会明显影响题目所问的量？
  \item 忽略它以后，主要规律是否变得清楚、可计算？
  \item 如果题目条件改变，这个模型还是否成立？
\end{enumerate}

\textbf{认知提醒：} 模型不是否定真实世界，而是先抓主要因素。解题时要能说出“为什么可以这样简化”，也要知道条件改变时模型可能失效。

\end{conceptbox}

\begin{conceptbox}{矢量与标量}

\textbf{直观理解：}
\begin{itemize}
  \item \textbf{标量}：只有大小的量。比如温度 $30^\circ\mathrm{C}$、质量 $5\ \mathrm{kg}$、路程 $100\ \mathrm{m}$。用一个数值和一个单位就能完整描述。
  \item \textbf{矢量}：既有大小又有方向的量。比如位移“向东 $100\ \mathrm{m}$”、速度“向南 $20\ \mathrm{m/s}$”、力“向上 $10\ \mathrm{N}$”。方向和大小一样重要。
\end{itemize}

\textbf{为什么要区分？} 两个标量直接加减就行：$5\ \mathrm{kg} + 3\ \mathrm{kg} = 8\ \mathrm{kg}$。但两个矢量相加不能简单加减——向东走 $3\ \mathrm{m}$ 再向北走 $4\ \mathrm{m}$，合位移大小不是 $7\ \mathrm{m}$，而是 $5\ \mathrm{m}$（勾股定理），方向为东偏北。方向不同，算法不同。

\textbf{严格定义：} 矢量是满足平行四边形定则（或三角形定则）的物理量，具有大小和方向。标量只具有大小，只遵循代数加减法则。

\textbf{平行四边形法则：} 把两个矢量平移到同一起点，以这两个矢量为邻边作平行四边形，从共同起点指向对角点的对角线就是合矢量。

\begin{center}
% TikZ 图：矢量加法示意（东3北4得5）
% Day 01：初高中物理的区别与学习方法
% 使用：% TikZ 图：矢量加法示意（东3北4得5）
% Day 01：初高中物理的区别与学习方法
% 使用：% TikZ 图：矢量加法示意（东3北4得5）
% Day 01：初高中物理的区别与学习方法
% 使用：\input{figures/tikz/vector-addition.tex}
\begin{tikzpicture}[scale=1.4,>=Latex]
  % 画箭头（加长让标注有空间）
  \draw[->,thick] (0,0) -- (3.2,0);
  \draw[->,thick] (3.2,0) -- (3.2,4.2);
  \draw[->,ultra thick,red] (0,0) -- (3.2,4.2);
  % 长度标注放在箭头正中间的上方/右侧，远离坐标轴交汇处
  \node[above,yshift=2pt] at (1.6,0) {$3\ \mathrm{m}$};
  \node[right,xshift=4pt] at (3.2,2.1) {$4\ \mathrm{m}$};
  \node[above left,yshift=4pt,xshift=-2pt] at (1.6,2.1) {$5\ \mathrm{m}$};
  % 方向标注放在箭头末端之外
  \node[below] at (3.2,-0.15) {东};
  \node[right] at (3.5,4.2) {北};
  % 起点
  \node[below left,yshift=-2pt] at (0,0) {起点};
\end{tikzpicture}

\begin{tikzpicture}[scale=1.4,>=Latex]
  % 画箭头（加长让标注有空间）
  \draw[->,thick] (0,0) -- (3.2,0);
  \draw[->,thick] (3.2,0) -- (3.2,4.2);
  \draw[->,ultra thick,red] (0,0) -- (3.2,4.2);
  % 长度标注放在箭头正中间的上方/右侧，远离坐标轴交汇处
  \node[above,yshift=2pt] at (1.6,0) {$3\ \mathrm{m}$};
  \node[right,xshift=4pt] at (3.2,2.1) {$4\ \mathrm{m}$};
  \node[above left,yshift=4pt,xshift=-2pt] at (1.6,2.1) {$5\ \mathrm{m}$};
  % 方向标注放在箭头末端之外
  \node[below] at (3.2,-0.15) {东};
  \node[right] at (3.5,4.2) {北};
  % 起点
  \node[below left,yshift=-2pt] at (0,0) {起点};
\end{tikzpicture}

\begin{tikzpicture}[scale=1.4,>=Latex]
  % 画箭头（加长让标注有空间）
  \draw[->,thick] (0,0) -- (3.2,0);
  \draw[->,thick] (3.2,0) -- (3.2,4.2);
  \draw[->,ultra thick,red] (0,0) -- (3.2,4.2);
  % 长度标注放在箭头正中间的上方/右侧，远离坐标轴交汇处
  \node[above,yshift=2pt] at (1.6,0) {$3\ \mathrm{m}$};
  \node[right,xshift=4pt] at (3.2,2.1) {$4\ \mathrm{m}$};
  \node[above left,yshift=4pt,xshift=-2pt] at (1.6,2.1) {$5\ \mathrm{m}$};
  % 方向标注放在箭头末端之外
  \node[below] at (3.2,-0.15) {东};
  \node[right] at (3.5,4.2) {北};
  % 起点
  \node[below left,yshift=-2pt] at (0,0) {起点};
\end{tikzpicture}


\figcaption{矢量加法的平行四边形法则：向东 $3\ \mathrm{m}$ 与向北 $4\ \mathrm{m}$ 的合矢量大小为 $5\ \mathrm{m}$（粗对角线）}\label{fig:vector-addition}
\end{center}

\vspace{0.3em}
常见矢量：位移、速度、加速度、力、动量、电场强度、磁感应强度\\
常见标量：质量、时间、温度、路程、速率、功、能量、电势

\end{conceptbox}

\begin{conceptbox}{位移与路程：先建立一个入口}

Day 02 会系统学习位置、位移、路程和速度；这里先把后面反复要用的一组区别说清。

\begin{itemize}
  \item \textbf{位移}描述从初位置到末位置的变化，是矢量；既要说大小，也要说方向。
  \item \textbf{路程}描述实际走过的路径长度，是标量；只记长度，不记方向。
  \item 因而总有 $|\Delta\vect{x}|\leq s$。绕一圈回到出发点时，位移为零，路程却不为零。
\end{itemize}

例如，先向东走 $3\ \mathrm{km}$，再向南走 $4\ \mathrm{km}$：路程是 $7\ \mathrm{km}$；位移是从起点直指终点的合矢量，大小为 $5\ \mathrm{km}$。这里不能把两个长度直接相加得到位移大小，原因正是位移有方向。Day 02 将用平行四边形法则和三角形法则画出这一过程。

\end{conceptbox}

\begin{conceptbox}{国际单位制（SI）}

物理公式里的数字不是凭空来的。同一个物理量用不同单位，数值完全不同。为了统一，全世界使用\textbf{国际单位制（SI）}。

国际单位制有七个基本单位。高中阶段应重点会使用米、千克、秒、安培、开尔文和摩尔；坎德拉了解即可。
\[
\text{米 (m)、千克 (kg)、秒 (s)、安培 (A)、开尔文 (K)、摩尔 (mol)、坎德拉 (cd)}
\]

其中，摩尔 $\mathrm{mol}$ 是物质的量的单位，在化学、热学和气体问题中都会用到，也属于高中阶段需要会使用的基本单位。

从基本单位可以导出所有其他单位：
\[
\text{速度：}\ \mathrm{m/s},\qquad
\text{加速度：}\ \mathrm{m/s^2},\qquad
\text{力：}\ \mathrm{kg\cdot m/s^2 = N}
\]

\textbf{弧度与角度：} 圆周运动中常用弧度 $\mathrm{rad}$ 描述角度。弧度的定义是
\[
\theta=\frac{s}{r},
\]
即“弧长除以半径”。因此弧度对应两个长度之比，是无量纲量；写成 $\mathrm{rad}$ 是为了提醒这是角度。后面使用 $s=r\theta$、$\omega=\dfrac{\Delta\theta}{\Delta t}$、三角函数小角度近似时，$\theta$ 都要用弧度而不是角度制。

\textbf{常见陷阱：} 本讲义默认在代入公式前统一为 SI 单位，例如把 $\mathrm{km/h}$ 换算成 $\mathrm{m/s}$。协调的非 SI 单位也能用于计算，但初学阶段统一使用 SI 最稳妥。
\[
1\ \mathrm{km/h} = \frac{1000\ \mathrm{m}}{3600\ \mathrm{s}} = \frac{1}{3.6}\ \mathrm{m/s}
\]

质量单位中的 $\mathrm{g}$ 表示“克”，而 SI 基本单位是 $\mathrm{kg}$。因为 $1\ \mathrm{kg}=1000\ \mathrm{g}$，所以把克换成千克要除以 $1000$：
\[
1\ \mathrm{g}=10^{-3}\ \mathrm{kg}
\]
\textbf{符号检查：} 看到 $\mathrm{g}$ 先判断它是不是单位“克”；看到斜体 $g$ 再按重力加速度处理。两者不是同一个物理量，不能直接混用。

\vspace{0.3em}
\textbf{测量意识：} 物理量要写数值、单位和合理的有效数字。同一记录规则下，$2.00\ \mathrm{m}$ 比 $2.0\ \mathrm{m}$ 多保留一位，表示更细的报告分辨率，但不代表一定更接近真值。测量总有误差，不能把结果看成绝对精确值。

\end{conceptbox}

\section*{四、公式与适用条件}

\begin{samepage}
\textbf{建模与单位使用前检查：}
\nopagebreak[4]
\begin{enumerate}
  \item 先看题目问的是“描述现象”，还是“计算某个物理量”。
  \item 判断研究对象能否简化为质点、刚体、光滑面等理想模型。
  \item 区分题中量是矢量还是标量，矢量必须说明方向或正方向。
  \item 代入任何公式前，先把单位统一到国际单位制，并用量纲粗略检查。
\end{enumerate}
\end{samepage}

本节没有需要记忆的力学公式，但有一个贯穿全书的重要工具——\textbf{量纲分析}。

\begin{conceptbox}{量纲分析初步}

\textbf{概念：} 量纲表示一个物理量由哪些基本量组合而成，反映的是“物理量的类型”，不是某一个具体单位。比如 $\mathrm{m}$ 和 $\mathrm{km}$ 都是长度单位，量纲都是 $[L]$；$\mathrm{m/s}$ 和 $\mathrm{km/h}$ 都是速度单位，量纲都是 $[L][T]^{-1}$。

\begin{itemize}
  \item 长度量纲：$[L]$
  \item 质量量纲：$[M]$
  \item 时间量纲：$[T]$
\end{itemize}

速度和加速度的量纲可以直接从定义看出来：
\[
v=\frac{\Delta x}{\Delta t},\qquad a=\frac{\Delta v}{\Delta t}
\]
所以：
\[
[\text{速度}] = [L][T]^{-1},\qquad
[\text{加速度}] = [L][T]^{-2},\qquad
[\text{力}] = [M][L][T]^{-2}
\]

\textbf{用途：} 检验公式是否合理。一个公式等号两边的量纲必须相同。

\textbf{举例：} 对匀变速直线运动的候选关系
$\Delta x = v_0 t + \frac{1}{2} a t^2$ 作量纲检查。这里暂不判断它的适用条件和数值系数，只检查各项能否相加：
\begin{itemize}
  \item 左边位移 $\Delta x$ 的量纲是 $[L]$
  \item 右边 $v_0 t$ 的量纲：$[L][T]^{-1} \times [T] = [L]$ \quad (量纲一致)
  \item 右边 $\frac{1}{2} a t^2$ 的量纲：$[L][T]^{-2} \times [T]^2 = [L]$ \quad (量纲一致)
  \item 左右量纲一致，说明该关系通过了量纲检查，但还不能据此证明公式正确
\end{itemize}

\textbf{边界提醒：} 量纲分析只能判断公式是否可能正确，不能证明它绝对正确（系数 $\frac{1}{2}$ 是否正确就需要通过推导或实验来确认）。

\end{conceptbox}

\section*{五、典型模型}

\textbf{质点模型}——高中物理最基本的模型，贯穿整个运动学和力学。

\begin{conceptbox}{质点}
质点是一个\textbf{理想物理模型}：把物体简化为一个有质量的点，忽略大小和形状，只保留质量和位置等关键信息。

\begin{itemize}
  \item \textbf{什么时候能用？} 物体的形状和大小对研究问题没有影响，或影响可以忽略时。
  \item \textbf{典型例子：}
    \begin{itemize}
      \item 研究地球绕太阳公转——地球可视为质点（因为公转半径远大于地球半径）
      \item 研究地球自转——地球不能视为质点（自转跟形状直接相关）
      \item 研究火车从北京到上海——火车可视为质点（距离远大于火车长度）
      \item 研究火车完全通过一座桥所需的时间——若桥长和火车长度可比，火车不能视为质点
    \end{itemize}
  \item \textbf{核心思想：} 同一个物体，在不同问题中能不能看成质点，取决于题目所关心的问题，而不取决于物体本身的大小。
\end{itemize}
\end{conceptbox}

\nopagebreak[4]
\section*{六、例题精讲}

\begin{examplebox}{判断物体能否视为质点}

\textbf{题意分析：} 判断以下四种情况中，研究对象能否视为质点，并说明理由。

\begin{enumerate}
  \item 研究跳水运动员在空中翻转的动作。
  \item 研究运动员质心从起跳到入水所经过的轨迹长度。
  \item 研究一列火车通过 $1000\ \mathrm{m}$ 长的隧道所需时间（火车长 $200\ \mathrm{m}$）。
  \item 研究地球绕太阳公转一周的时间。
\end{enumerate}

\textbf{物理模型：} 质点模型。判断依据是：物体的形状和大小对研究的问题是否产生不可忽略的影响。

\textbf{解题过程：}
\begin{enumerate}
  \item \textbf{不能。} 翻转动作依赖于身体的姿态和转动，必须知道身体各部分的相对位置，不能把身体简化为一个点。
  \item \textbf{能。} 题目已经明确研究质心轨迹；若运动员的大小和姿态不影响质心轨迹长度的计算，就可用一个质点的位置变化来描述。
  \item \textbf{不能。} 火车长度 $200\ \mathrm{m}$ 与隧道长度 $1000\ \mathrm{m}$ 可比，计算通过时间需要考虑火车长度。如果隧道长 $100\ \mathrm{km}$，火车可视为质点。
  \item \textbf{能。} 公转半径约 $1.5\times 10^{11}\ \mathrm{m}$，地球半径约 $6.4\times 10^{6}\ \mathrm{m}$，前者是后者的约 $23400$ 倍，在本题要求的精度下可以忽略地球大小。
\end{enumerate}

\textbf{方法总结：} 判断质点时先问“形状、大小、转动会不会影响答案”。若不影响，就可简化为质点；若影响，就必须保留物体的形状或内部结构。物体本身大小不是唯一标准。

\end{examplebox}

\section*{七、变式训练}

\begin{enumerate}
  \item 判断：研究乒乓球在发球时的旋转轨迹，乒乓球能否视为质点？研究乒乓球落台后的运动轨迹（忽略旋转），乒乓球能否视为质点？分别说明理由。
  \item 一个人向东走了 $3\ \mathrm{km}$，再向南走了 $4\ \mathrm{km}$。求他的总路程和位移大小。这个例子说明路程和位移有什么区别？
  \begin{center}
  \begin{tikzpicture}[scale=0.82,>=Stealth]
  \coordinate (O) at (0,0);
  \coordinate (A) at (3.4,0);
  \coordinate (B) at (3.4,-2.35);
  \draw[->,thick] (O) -- (A) node[midway,above] {$3\ \mathrm{km}$};
  \draw[->,thick] (A) -- (B) node[pos=0.38,right=6pt] {$4\ \mathrm{km}$};
  \draw[->,blue,thick] (O) -- (B) node[pos=0.43,below left] {位移};
  \node[left] at (O) {起点};
  \node[right] at (B) {终点};
  \node[below] at (1.7,-0.06) {向东};
  \node[right] at (4.12,-1.55) {向南};
\end{tikzpicture}

  \end{center}
  \item 用量纲分析判断：公式 $v^2 = v_0^2 + 2ax$ 是否量纲正确？（提示：速度 $v=\Delta x/\Delta t$，加速度 $a=\Delta v/\Delta t$，$x$ 是位移）
\end{enumerate}

\begin{learningbox}{变式训练参考结果}
\begin{enumerate}[itemsep=0.45em,topsep=0.3em,parsep=0pt]
  \item 研究乒乓球的旋转时，球的大小和转动不能忽略，不能视为质点；研究其落台后的运动轨迹且忽略旋转时，球的大小相对运动范围可忽略，可以视为质点。
  \item 总路程为 $3+4=7\ \mathrm{km}$；位移大小为
  $\sqrt{3^2+4^2}=5\ \mathrm{km}$，方向为东偏南。路程是实际运动路径的长度，是标量；位移由初位置指向末位置，是矢量。
  \item $[v^2]=\mathrm{L^2T^{-2}}$，$[v_0^2]=\mathrm{L^2T^{-2}}$，
  $[ax]=\mathrm{LT^{-2}}\cdot\mathrm L=\mathrm{L^2T^{-2}}$，等式各项量纲相同，公式量纲正确。
\end{enumerate}
\end{learningbox}

\section*{八、竞赛入门补充}

\begin{competitionbox}
\textbf{量纲分析的进阶用法——推导物理关系}

如果已知一个物理量可能依赖哪些因素，可以用量纲分析快速确定它们之间的关系。这是竞赛中常用的初步判断方法。

\textbf{例题：} 假设单摆的周期 $T$ 只可能与摆长 $L$、重力加速度 $g$、摆球质量 $m$ 和最大摆角 $\theta_0$ 有关。用量纲分析确定周期表达式中能够由量纲决定的部分。

\textbf{思路：}
\begin{center}
\begin{minipage}[c]{0.42\linewidth}
\centering
% TikZ 图：Day 01 单摆量纲分析变量
\begin{tikzpicture}[scale=0.72,>=Latex,line cap=round,line join=round,
  every node/.style={font=\small}]
  \coordinate (O) at (0,0);
  \coordinate (B) at (-1.45,-2.55);

  \draw[thick] (-1.15,0) -- (1.15,0);
  \foreach \x in {-1.0,-0.7,...,1.0}
    \draw[gray!65] (\x,0) -- ++(-0.18,0.18);
  \fill (O) circle (1.5pt);

  \draw[dashed,gray!75] (O) -- (0,-2.85);
  \draw[very thick] (O) -- (B);
  \filldraw[fill=white,thick] (B) circle (0.20);
  \node[left=4pt] at (B) {$m$};

  \draw[->,blue!75!black,thick] (B) -- ++(0,-0.92)
    node[below] {$\vect g$};
  \draw[<->] ($(O)+(-0.18,0.08)$) -- ($(B)+(-0.18,0.08)$)
    node[midway,left=3pt] {$L$};
  \draw (0,-0.55) arc[start angle=-90,end angle=-119.6,radius=0.55];
  \node at (-0.34,-0.78) {$\theta_0$};
\end{tikzpicture}

\end{minipage}\hfill
\begin{minipage}[c]{0.50\linewidth}
\[
\begin{aligned}
[T]&=[T], & [L]&=[L],\\
[g]&=[L][T]^{-2}, & [m]&=[M].
\end{aligned}
\]
最大摆角 $\theta_0$ 用弧度表示，是无量纲量。
\end{minipage}

\figcaption{单摆量纲分析中的变量}\label{fig:day01-dimensional-pendulum}
\end{center}

因此 $\theta_0$ 可以通过某个无量纲函数 $f(\theta_0)$ 影响周期，而这种函数关系无法由量纲确定。周期 $T$ 的量纲中不包含质量 $[M]$；在题设给出的变量里，只有 $m$ 带有质量量纲，且没有其他量可以与它抵消，所以 $m$ 的指数只能为 $0$。设
\[
T=L^a g^b f(\theta_0),
\]
则
\[
[T] = [L]^a \cdot ([L][T]^{-2})^b = [L]^{a+b}[T]^{-2b}
\]
对比两边：$a + b = 0$（$L$ 指数），$-2b = 1$（$T$ 指数）
解得 $b = -\frac12$，$a = \frac12$，所以量纲分析只能得到：
\[
T=\sqrt{\frac{L}{g}}\,f(\theta_0)
\]

在小摆角近似下，进一步的动力学推导给出
$f(\theta_0)\approx2\pi$，于是
$T\approx2\pi\sqrt{L/g}$；摆角较大时，$f(\theta_0)$ 会随最大摆角变化。量纲分析既不能给出常数 $2\pi$，也不能确定函数 $f$ 的具体形式。

\textbf{竞赛提示：} 物理竞赛中，如果遇到陌生的公式判断，先用量纲分析筛掉明显错误的选项，可以节省大量时间。
\end{competitionbox}

\vfill
\pagebreak
\section*{九、易错点}

\begin{mistakebox}
\begin{enumerate}
  \item \textbf{把路程当位移：} 先看题目问“实际走过多长”还是“从起点到终点的变化”。位移是矢量，看起点到终点的直线距离和方向；路程是标量，看实际路径总长。
  \item \textbf{认为“小物体就是质点”：} 先问形状、大小、转动是否影响本题。若不影响，可看作质点；若题目研究转动、形变或内部结构，就不能这样简化。
  \item \textbf{忽略单位换算：} 列式前先统一 SI 单位。看到 $\mathrm{km/h}$ 除以 3.6，看到 $\mathrm{cm}$ 除以 100，看到质量单位 $\mathrm{g}$（克）除以 1000，换成 $\mathrm{kg}$；不要把“克”的符号 $\mathrm{g}$ 与重力加速度 $g$ 混淆。
  \item \textbf{公式中的常数系数不能用量纲判断：} 量纲分析先筛掉结构明显错误的式子；像 $\frac{1}{2}at^2$ 中的 $\frac12$，还要靠推导或实验确定。
\end{enumerate}
\end{mistakebox}

\section*{十、今日小结}

\begin{enumerate}
  \item 高中物理更强调以模型为起点，把定性判断、定量关系、方向和适用条件组织成完整推理。
  \item 物理模型是对现实的合理简化。质点是最基本的模型——忽略大小和形状，保留质量、位置及研究问题所需的运动信息。
  \item 矢量和标量的根本区别：矢量有方向，遵循平行四边形定则；标量只有大小，直接代数加减。
  \item 国际单位制（SI）是高中物理定量计算的共同语言；本讲义默认在使用公式前统一为 SI 单位。
  \item 量纲分析是检验公式是否合理的有力工具，但无法确定公式中的常数系数。
\end{enumerate}

\textbf{一分钟复盘：} 不看正文，画一个三列表：模型、方向、单位。模型栏写“什么时候可看成质点”；方向栏写“速度为什么要标方向”；单位栏写“代公式前怎样统一 SI 单位”。三栏都能填出例子，后面的运动学学习就有了清晰入口。
